\centerline{\bf \Huge Финальная МатДрака }

\problem{\textbf{1!}} Чётно или нечётно число  23 + 26 + 29 + ... + $X$? Если 23 первое слагаемое, 26 второе, 29 третье, а $X$ две тысячи двадцать пятое слагаемое.


\problem{\textbf{2}} Замените в равенстве   ПИРОГ = КУСОК + КУСОК + КУСОК + ... + КУСОК   одинаковые буквы одинаковыми цифрами, а разные – разными так, чтобы равенство было верным, а количество "кусков пирога" было бы наибольшим из возможных. Приведите хотя бы два примера!


\problem{\textbf{3}} Все натуральные числа, сумма цифр в записи которых делится
на 5, выписывают в порядке возрастания: 5, 14, 19, 23, 28, 32, . . . Чему равна самая маленькая
положительная разность между соседними числами в этом ряду? Приведите пример таких чисел!


\problem{\textbf{4!}} Можно ли составить магический квадрат из первых 49 простых чисел?
Магический квадрат – это квадратная таблица, заполненная числами, в которой суммы чисел во всех строках и столбцах равны.


\problem{\textbf{5}} Найдите количество натуральных чисел от 1 до 100,
имеющих ровно четыре натуральных делителя, не менее чем три из которых не превосходят 10.


\problem{\textbf{6}} Из 1000 ребят, обучающихся на Летних сборах ЭлМШ, 334 смотрели “Гарри Поттера”, 308 смотрели “Звездные
войны”, 293 смотрели “Властелина колец”, 266 смотрели “Хроники Нарнии”. “Звездные войны” и “Гарри Поттера” смотрели 79
ребят. “Звездные войны” и “Властелина колец” смотрели 67 ребят. “Звездные войны” и “Хроники Нарнии” смотрели 59 ребят.
“Гарри Поттера” и “Властелина колец” смотрели 32 человека. “Гарри Поттера” и “Хроники Нарнии” смотрели 125 ребят.
“Властелина колец” и “Хроники Нарнии” смотрели 110 ребят. Все 4 фильма смотрели 26 ребят. При этом если кто-то не видел
один из этих 4 фильмов, он обязательно не видел и ещё один из них. Сколько ребят не смотрели ни один из этих фильмов?


\problem{\textbf{7}} Назовём трёхзначное число интересным, если хотя бы одна его
цифра делится на 3. Какое наибольшее количество подряд идущих интересных чисел может
быть? Приведите пример таких чисел!



\problem{\textbf{8}}  Сколькими способами цифры от 1 до 9 можно разбить на
несколько (больше одной) групп так, чтобы суммы цифр во всех группах были равны друг
другу? (Разбиения, отличающиеся только перестановкой групп, считаются одинаковыми.)

\newpage
\hspace{6mm}

\problem{\textbf{9}} Произведение пяти чисел не равно нулю. Каждое из этих чисел уменьшили на единицу, при этом их произведение не изменилось. Приведите пример таких чисел!


\problem{\textbf{10}} Назовём натуральное число «примечательным»,
если все его цифры попарно различны и их сумма равна 18. Найдите сумму примечательных
чисел, не превосходящих 950.

\problem{\textbf{11}} Сколько существует натуральных чисел, меньших тысячи, которые не делятся
ни на 5, ни на 7?




\problem{\textbf{12}} В кошельке у купца Ганса лежат 20 серебряных
монет по 2 кроны, 15 серебряных монет по 3 кроны и 3 золотых дуката (1 дукат равен 5 крон).
Сколькими способами Ганс может уплатить сумму в 10 дукатов? Монеты одного достоинства
неразличимы.


\problem{\textbf{13}} Сумма десяти натуральных чисел равна 1001. Какое наибольшее
значение может принимать НОД (наибольший общий делитель) этих чисел?



\problem{\textbf{14}} Сколько существует пар различных натуральных чисел таких, что их
произведение на 19999 больше их суммы?



\problem{\textbf{15}} Каждая буква в словах ЭХ и МОРОЗ соответствует какой-то цифре, причём одинаковым цифрам соответствуют одинаковые буквы, а разным – разные.

Известно, что  Э$\times$Х = M$\times$О$\times$Р$\times$О$\times$З,  а  Э + Х = М + О + Р + О + З. 

Чему равно  Э$\times$Х + M$\times$О$\times$Р$\times$О$\times$З?

